% Candidate only; not inserted into the current eight-page manuscript.
% New native-CPU protocol; E10 is a conservative lift diagnostic, not Table III's SR10.
\begin{table}[t]
\centering\footnotesize
\caption{Recovery in novel obstacle arrangements in simulation.}
\label{tab:scene_generalization}
\begin{tabular}{lrrr}
\toprule
Scene & $N$ & Quiet hold$_{10}$ (\%) & Lift-clear hold$_{10}$ (\%)\\
\midrule
Crossed timber & 120 & 86.7 & 77.5\\
Ladder and boards & 120 & 89.2 & 76.7\\
Hollow container & 120 & 87.5 & 81.7\\
\midrule
Fixed overhead enclosure & 120 & 65.0 & 64.2\\
\bottomrule
\end{tabular}
\par\vspace{2pt}
\begin{minipage}{.98\columnwidth}\scriptsize
Each scene includes 30 trials per initial orientation and 75\% natural/25\% procedural low states. Trials last 20\,s; both columns require a continuous 10\,s hold. The latter additionally requires the rigid vertical-lift collision probe to remain clear. This diagnostic may reject standing beside residual objects and is not equivalent to articulated escape feasibility. Source poses recur across layouts; these are single-policy simulation results, not hardware statistics or multi-seed estimates.
\end{minipage}
\end{table}
